\begingroup
\renewcommand{\arraystretch}{1.45}
\fancyhead[L]{}
\fancyhead[C]{\small Das wunderbare Pentagramm}
\par\medskip\Needspace{5\baselineskip}\pdfbookmark[1]{Art. 7}{pentagramma-art-7}\centerline{[7.]}\nobreak\smallskip
1843. April 20. Die exzentrischen Anomalien $\varphi,\varphi',\varphi'',\varphi''',\varphi''''$ der Punkte
$A,B,C,D,E$ sind durch die Gleichungen verbunden ($G$ als negativ betrachtet)
\[
\begin{gathered} \frac{\sin\frac12(\varphi'''+\varphi'')}{\cos\frac12(\varphi'''-\varphi'')} =\frac{G}{G''}\mathbin{.}\sin\varphi,\\[9pt]\quad{}\qquad \frac{\cos\frac12(\varphi'''+\varphi'')}{\operatorname{cso}\frac12(\varphi'''-\varphi'')} =\frac{G}{G''}\mathbin{.}\cos\varphi\\[9pt] \frac{\sin\frac12(\varphi'+\varphi'''')}{\cos\frac12(\varphi'-\varphi'''')} =\surd\frac{G(G-1)}{G''(G''-1)}\mathbin{.}\sin\varphi \\[9pt]\quad{}=\frac{G(2G-1}{G''(2G''-1)}\sin\varphi\\[9pt] \frac{\cos\frac12(\varphi'+\varphi'''')}{\cos\frac12(\varphi'-\varphi'''')} =\surd\frac{G(G-1)}{G'(G'-1)}\mathbin{.}\cos\varphi \\[9pt]\quad{}=\frac{G(2G-1)}{G'(2G'-1)}\cos\varphi \end{gathered}
\]

\Needspace{30\baselineskip}
Die Beziehungen zwischen den Winkeln $\varphi^0,\varphi',\varphi''$ sind am einfachsten auf
folgende Art darzustellen, $\surd\frac{G'}{G'-1}=\xi$, $\surd\frac{G''}{G''-1}=\eta$ gesetzt, wird $\xi\eta=$
\[
\begin{gathered}\frac{\operatorname{tang}\frac12(\varphi'-\varphi'''')}{\operatorname{tang}\frac12(\varphi'''-\varphi'')} \\[9pt]\quad{}=\frac{\operatorname{tang}\frac12(\varphi''-\varphi^0)}{\operatorname{tang}\frac12(\varphi''''-\varphi''')} \\[9pt]\quad{}=\frac{\operatorname{tang}\frac12(\varphi'''-\varphi')}{\operatorname{tang}\frac12(\varphi^0-\varphi'''')} \\[9pt]\quad{}=\frac{\operatorname{tang}\frac12(\varphi''''-\varphi'')}{\operatorname{tang}\frac12(\varphi'-\varphi^0)} \\[9pt]\quad{}=\frac{\operatorname{tang}\frac12(\varphi^0-\varphi''')}{\operatorname{tang}\frac12(\varphi''-\varphi')}\end{gathered}
\]
für $\frac{\xi\eta+1}{\xi\eta-1}$ gibt es einen ähnlichen Ausdruck, der sich hieraus leicht ableiten lässt.



\begin{landscape}
\SourceLeaf{499}

In Zahlen
\par\nopagebreak% portrait-table start
\begingroup\small\relax\setlength{\arraycolsep}{2pt}\[
\resizebox{.98\linewidth}{!}{$\displaystyle \begin{array}{r@{\;}c@{\;}rrr@{\quad}rrr@{\quad}rrr@{\quad}rrr} &&&&&&&&&&&\log\operatorname{tang}&\log\operatorname{tang}&\log\Delta\\[9pt] \varphi^0&=&50^0&29'&20''&80^0&54'&55''&49^0&13'&4''&0.79616&0.06418&9.81007\\[9pt] \varphi'&=&92&56&38&55&49&27&15&16&12.5&0.16814&9.43617&9.182\\[9pt] \varphi''&=&162&8&14&83&48&52&59&41&16.5&0.96505&0.23313&9.97901\\[9pt] \varphi'''&=&260&34&22&64&29&16.5&21&13&39&0.32127&9.58931&9.33515\\[9pt] \varphi''''&=&291&6&47&74&57&29&34&35&48&0.57068&9.83871&9. \end{array}$}
\]\endgroup\par% portrait-table end

\hfill $0.73196=\log\xi\eta$\hspace*{.12\textwidth}\par


\end{landscape}
\par\medskip\Needspace{5\baselineskip}\pdfbookmark[1]{Art. 8}{pentagramma-art-8}\centerline{[8.]}\nobreak\smallskip
Die $\varphi,\varphi',\varphi'',\varphi''',\varphi''''$ sind nichts anderes als die Amplituden zu fünf transzendenten Argumenten, welche um $\frac45K$ zunehmen (in der Bedeutung von \textsc{Jacobi}
p.\,31) und wo der Modul $k$
\[
\begin{gathered}=\surd\frac{\frac{1}{G'G'}-\frac{1}{G''G''}}{\frac{1}{G'G'}-\frac{1}{GG}} =\sin\mu,\\[9pt]\quad{}\qquad \cos\mu=\surd\frac{\frac{1}{G''G''}-\frac{1}{GG}}{\frac{1}{G'G'}-\frac{1}{GG}}\end{gathered}
\]
Das transzendente Argument selbst, unbestimmt genommen, ist
\[
=\int\frac{x\,\mathrm{d}y-y\,\mathrm{d}x}{\surd(xx+yy)}\mathbin{.} \surd\frac{\frac{1}{G'G'}-\frac{1}{GG}} {\left(\frac{1}{G'G'}-\frac{1}{GG}\right)xx+\left(\frac{1}{G''G''}-\frac{1}{GG}\right)yy}
\]

\Needspace{23\baselineskip}
$\Delta$ wird in der Bezeichnung von \textsc{Jacobi} verwendet, sodass $\Delta\varphi=\surd(1-kk\sin\varphi^2)$,
es sind
\par\nopagebreak% portrait-table start
\begingroup\small\setlength{\arraycolsep}{2pt}\[
\resizebox{.98\linewidth}{!}{$\displaystyle \begin{array}{c|c|c|c|c} & & &\text{propositional}&\\[9pt] \text{die drei Größen}&\text{ebenso}&&\text{den Zahlen}&\text{oder}\\[9pt] \operatorname{tang}\frac12(\varphi'-\varphi'''')& \operatorname{tang}\frac12(\varphi''-\varphi^0)&\text{usw.}& G(2G-1)\surd\left(\frac{1}{G'G'}-\frac{1}{GG}\right)& \operatorname{tang}\operatorname{am}\frac45K\\[9pt] \operatorname{tang}\frac12(\varphi'''-\varphi'')& \operatorname{tang}\frac12(\varphi''''-\varphi''')&& -G\surd\left(\frac{1}{G'G'}-\frac{1}{GG}\right)& \operatorname{tang}\operatorname{am}\frac25K\\[9pt] \Delta\varphi^0&\Delta\varphi'&&1&1 \end{array}$}
\]\endgroup\par% portrait-table end

\par\nopagebreak\begin{center}\rule{.25\textwidth}{.4pt}\end{center}
\clearpage
\endgroup
